\begin{center}
    \textbf{ВЫВОДЫ}
\end{center}
\addcontentsline{toc}{chapter}{ВЫВОДЫ}

Оба семейства операционных систем Windows и UNIX являются системами разделения времени с динамическими приоритетами и вытеснением, поэтому обработчик прерывания от системного таймера выполняет в этих системах схожие функции:

\begin{enumerate} 
	\item выполняют декремент кванта;
	\item инкременитирует счетчики времени и декрементирует счетчики таймеров, будильников реального времени, счетчики времени отложенных задач;
	\item инициализирует отложенные действия, относящиеся к работе планировщика, такие как пересчёт приоритетов.

\end{enumerate} 

Пересчет динамических приоритетов осуществляется только для пользовательских потоков для того, чтобы избежать их бесконечного откладывания. В ОС семейства UNIX приоритеты вычисляются по формуле, в которой учитываются время простоя в очереди готовых потоков и полученное процессорное время. В ОС семейства Windows пересчет приоритетов осуществляется путем сканирования очереди готовых потоков.
